\begin{UseCase}{CU-49}{Registrar composición del núcleo familiar}{
	Permite al Trabajador Social registrar el número de personas en el hogar, quiénes conviven con el NNA y si algún hermano vive fuera del hogar.
}
	\UCitem{Versión}{\color{Gray}1.0}
	\UCitem{Autor}{\color{Gray}Rivera Segura José Emiliano}
	\UCitem{Supervisa}{\color{Gray}Guerra Salinas Edgar Rafael}
	\UCitem{Actor}{\hyperlink{tTSocial}{Trabajador Social}}
	\UCitem{Propósito}{Documentar la composición del núcleo familiar para insumar el diagnóstico sociofamiliar y el PRD.}
	\UCitem{Entradas}{Número de personas en el hogar, descripción textual de convivientes, indicador de hermanos fuera del hogar, narración del entorno.}
	\UCitem{Origen}{Teclado.}
	\UCitem{Salidas}{Composición familiar registrada en el expediente.}
	\UCitem{Destino}{Pestaña Datos Familiares, subsección Núcleo familiar.}
	\UCitem{Precondiciones}{Usuario pertenece al equipo asignado.}
	\UCitem{Postcondiciones}{Datos del núcleo familiar almacenados.}
	\UCitem{Errores}{\begin{description}\item[E1:] Campo de número de personas vacío.\end{description}}
	\UCitem{Tipo}{Caso de uso primario}
	\UCitem{Observaciones}{Complementa el familiograma (\hyperlink{CU-27}{CU-27}).}
\end{UseCase}

\begin{UCtrayectoria}
	\UCpaso[\UCactor] Accede a la subsección \IUbutton{Núcleo Familiar}.
	\UCpaso Despliega el formulario de composición.
	\UCpaso[\UCactor] Ingresa el número de personas en el hogar y la descripción de quiénes conviven.
	\UCpaso[\UCactor] Indica si hay hermanos que viven fuera del hogar y captura la narración del entorno.
	\UCpaso[\UCactor] Presiona \IUbutton{Guardar composición familiar}.
	\UCpaso Valida campo de número de personas. \Trayref{A}.
	\UCpaso Almacena y muestra \textbf{MSG-09}.
	\UCpaso[] -- -- Fin del caso de uso.
\end{UCtrayectoria}

\begin{UCtrayectoriaA}{A}{Campo vacío}
	\UCpaso Muestra \textbf{MSG-04}. \UCpaso[\UCactor] Oprime \IUbutton{Aceptar}. \UCpaso Regresa al paso 3.
\end{UCtrayectoriaA}

%-------------------------------------- CU-50
